\section{Physical windows and the periodic gap}
\label{s:periodic-gap}

We pass from a uniform estimate on plane window marginals to a gap on
the physical torus.  All shell sets below consist of physical sites.
The finite coupling and Poincare tools used in this section are proved
in Appendix~\ref{a:finite}.  In particular, all heat-bath sums count
each indicated update at rate one.

\begin{definition}\label{d:shells}
For \(a\in\mathbb Z^2\), integer \(u\ge1\), put
\(C_u(a)=a+\{-u,\ldots,u\}^2\).
For \(\operatorname{dist}_\infty(y,C_u(a))=1\), let
\(s=u^{1/4}\), \(k_i=\lfloor is\rfloor\), \(i_1=\lfloor s\rfloor+1\).
For an integer \(l_0\ge1\) define
\[
 \begin{split}
 F_i&=\{x\in C_u(a):|x-y|_\infty>(i-1)s\},\\
 D_i&=\{x\in C_u(a):|x-y|_\infty=k_i\},\\
 A_i&=\{x\in D_i:\operatorname{dist}_\infty
                    (x,\mathbb Z^2\setminus C_u(a))\le l_0\},\qquad
 B_i=D_i\setminus A_i .
 \end{split}
\]
Let \(I_0\) consist of the integers \(2\le i\le i_1\) for which
some point of \(D_i\) is at distance strictly less than \(s/3\)
from a corner site of \(C_u(a)\).
All these are sets of physical sites.
\end{definition}

\begin{definition}[(W)]\label{d:W}
The family of plane frame specifications has \(\mathsf W(l_0,p_0,M_\star)\)
if \(l_0,M_\star\) are positive integers, \(0<p_0\le1\), and,
for every proper plane three-colouring, every translate \(C_u(a)\)
with integer \(u\ge M_\star\), every \(y\) as in Definition~\ref{d:shells},
every integer \(2\le i\le u^{1/4}\) outside \(I_0\), and every two
admissible frame-valued boundary conditions \(\eta,\xi\) for \(F_i\),
\[
 d_{\mathrm{TV}}\!\left(
       \mu^{+,\eta}_{\kappa,F_i,B_i},
       \mu^{+,\xi}_{\kappa,F_i,B_i}\right)\le1-p_0 .
\]
The subscript \(B_i\) denotes a marginal. Boundaries are arbitrary
within their site alphabets; they need not agree anywhere.
\end{definition}

For a positive specification with other finite site alphabets, the
phrase ``window bound (W)'' means the displayed marginal inequality
with that specification's laws in place of the frame laws, with the
same quantifiers over translates, exterior sites and permitted indices.
Consistency means that conditioning a volume law on its letters off a
subset gives the specification law on that subset with the induced
boundary.  Range one means that the latter law depends only on exterior
letters within \(\ell^\infty\) distance one of the subset.

\begin{lemma}\label{l:stopped}
Suppose a consistent positive plane specification with site-dependent
finite alphabets has range one in \(\ell^\infty\), one-site probabilities
at least \(\alpha\in(0,1]\), and the window bound (W) at every translate.
Put \(\delta=p_0\alpha^{16l_0}\) and
\(M_{\mathrm g}(l_0)=\max\{16,6(l_0+1)\}^4\).
For \(u\ge\max\{M_\star,M_{\mathrm g}(l_0)\}\), boundaries of \(C_u(a)\)
agreeing off one exterior site \(y\), and any integer
\(4\le h\le\lfloor u^{1/4}\rfloor\), there is a coupling of the
full laws such that
\[
 \mathbb P(n_{Q_h}\ne n'_{Q_h})\le(1-\delta)^{h-3},
 \qquad Q_h=\{x\in C_u(a):|x-y|_\infty>hu^{1/4}\}.
\]
If \(y\) is farther than range one, the full laws are equal.
\end{lemma}

\begin{proof}
Write \(C=C_u(a)\), \(\mu=\mu_C^\eta\) and
\(\mu'=\mu_C^\xi\) for the two volume laws.  If the exterior site
\(y\) is farther than distance one from \(C\), the range-one property
gives \(\mu=\mu'\); sampling the same configuration in both
components proves the conclusion.  Hence assume
\(\operatorname{dist}_\infty(y,C)=1\), and use the shells of
Definition~\ref{d:shells}.  The size assumption gives
\(s=u^{1/4}\ge\max\{16,6(l_0+1)\}\).
Lemma~\ref{l:shellgeom} supplies
\[
 |I_0|\le1,\qquad D_i\subset F_i,
 \qquad |A_i|\le8l_0\quad(i\notin I_0),
 \qquad 2\le i\le i_1.
\]
It also says that every site of \(C\setminus F_i\) within distance
one of \(F_i\) belongs to \(D_{i-1}\), and that \(y\) is not
within distance one of \(F_i\), for \(i\ge2\).
Since \(h+1\le\lfloor s\rfloor+1=i_1\), these geometric statements
apply to every surgery index used below.

We first record the consequence of finite energy needed for the
surgeries.  In any one of the positive finite-volume laws under
consideration, a one-site conditional probability after specifying
only some other letters is an average of fully specified one-site
conditional probabilities.  Thus it is at least \(\alpha\).
Order any finite set \(A\) and apply the finite chain rule to a
prescribed assignment \(t_A\), conditional on any fixed letters off
\(A\).  Every factor is at least \(\alpha\), so that
\[
 \mathbb P(n_A=t_A\mid\text{those fixed letters})
                      \ge\alpha^{|A|}.
\]
Positivity guarantees every conditioning denominator used in this
argument is nonzero.  For \(A=\varnothing\) the assertion reads
\(1\ge1\).

Start from the product coupling \(\mu\otimes\mu'\).  For the fixed
indices \(i=2,\ldots,h+1\), in increasing order, use the following
operation.  If \(i\in I_0\), leave the pair unchanged.  Otherwise,
for the current pair \((\omega,\omega')\), let \(p_i,p'_i\) be its
two specification laws on \(F_i\), with the respective complements
and original boundaries fixed.  If
\(\omega_{D_{i-1}}=\omega'_{D_{i-1}}\), these laws coincide:
all their exterior neighbours inside \(C\) belong to the agreeing
separator \(D_{i-1}\), and their exterior neighbours outside \(C\)
have identical boundary letters because none is \(y\).  In this
case draw one sample from the common law and use it in both components
on \(F_i\).

If the two configurations on \(D_{i-1}\) differ, maximally couple
the \(B_i\) marginals of \(p_i,p'_i\), using
Lemma~\ref{l:finitecoupling}.  Given the two sampled marginals,
complete them independently to \(F_i\) using their respective
conditional laws on \(F_i\setminus B_i\).  Finally resample \(A_i\)
independently in the two components, each with its conditional law
given the rest of its own configuration.  The initial completion has
the correct two marginals \(p_i,p'_i\).  The last resampling preserves
each of these laws, by Lemma~\ref{l:CE} and consistency.  Thus in
either branch the output on \(F_i\), for every fixed input pair,
has marginals \(p_i,p'_i\).  Lemma~\ref{l:surgery} shows that every
step preserves the original full marginals \(\mu,\mu'\).

Suppose now \(i\le h\) and \(i\notin I_0\).
Then \(i\le h\le\lfloor s\rfloor\), so (W) implies agreement
on \(B_i\) with probability at least \(p_0\) in the maximal-coupling
branch.  Fix an arbitrary assignment \(t_{A_i}\) using the alphabet
at each of its sites.  Given the entire result before the final
independent resamplings, the probability that both choose this
assignment is at least \(\alpha^{2|A_i|}\).
Those resamplings leave \(B_i\) fixed.  Therefore, conditional on
any history preceding step \(i\),
\[
 \mathbb P(\omega_{D_i}=\omega'_{D_i}
                     \text{ after step }i\mid\text{history})
       \ge p_0\alpha^{2|A_i|}
       \ge p_0\alpha^{16l_0}=\delta.
\]
In the identical-resampling branch the probability is one because
\(D_i\subset F_i\), so this bound holds in that branch as well.
At step \(h+1\), if \(h+1>s\), the coupling construction still
exists but no success bound from (W) is used.

Let
\[
 J=\{i\in\{2,\ldots,h\}:i\notin I_0,
                                      \ i+1\notin I_0\}.
\]
There are \(h-1\) possible trial indices, and the at most one
exceptional index excludes at most two of them.  Hence
\(|J|\ge h-3\ge1\).
If a trial \(i\in J\) produces agreement on \(D_i\), the
unskipped next step \(i+1\) uses identical resampling, giving agreement
on \(F_{i+1}\).  This agreement persists.  Indeed, for a later index
\(j\ge i+2\), \(F_j\subset F_{i+1}\), and
\[
 \lfloor(j-1)s\rfloor\ge(j-1)s-1
                      \ge(i+1)s-1>is,
\]
so \(D_{j-1}\subset F_{i+1}\).
A skipped step changes nothing, and any unskipped step therefore
again uses identical resampling on \(F_j\), preserving agreement
everywhere in \(F_{i+1}\).  Since
\(Q_h=F_{h+1}\subset F_{i+1}\), one successful trial in \(J\)
ensures final agreement on \(Q_h\).

To bound the chance of no success, list \(J=\{j_1<\cdots<j_r\}\)
and let \(H_i\) denote the full history through step \(i\).
The event that the first \(k-1\) trials fail is measurable with
respect to \(H_{j_k-1}\).  The displayed conditional success bound
therefore yields, successively,
\[
 \mathbb P(\text{first }k\text{ trials fail})
       \le(1-\delta)
           \mathbb P(\text{first }k-1\text{ trials fail}).
\]
Induction gives probability at most \((1-\delta)^r\), which is
at most \((1-\delta)^{h-3}\).  This proves the required bound for
the final coupling of the full laws.  If \(\delta=1\), the first
trial fails with probability zero, so the same argument gives zero.
An empty \(B_i\) has its unique empty configuration and is coupled
with agreement probability one; an empty \(A_i\) needs no resampling.
Empty volumes occurring in a step and singleton alphabets are covered
by the same formulas.  No comparison of letters at distinct sites
has been made.
\end{proof}

\begin{corollary}\label{p:surgery}
Under the hypotheses and size threshold of Lemma~\ref{l:stopped},
the same two laws have marginal total variation at most
\((1-\delta)^{u^{1/4}-4}\) on
\(\{x\in C_u(a):|x-y|_\infty>u^{1/2}\}\).
When \(\delta=1\), this upper bound is zero.
\end{corollary}

\begin{proof}
If the changed exterior site is farther than range one, the full laws
are equal by Lemma~\ref{l:stopped}.  Otherwise set
\(s=u^{1/4}\) and \(h=\lfloor s\rfloor\).
The size threshold gives \(s\ge16\), so \(h\ge4\).
Since \(hs\le s^2=u^{1/2}\), the target set of this corollary
is contained in \(Q_h\).  Lemmas~\ref{l:stopped}
and~\ref{l:finitecoupling} thus bound its marginal total variation by
\((1-\delta)^{h-3}\).
For \(0<\delta<1\), the inequality
\(\lfloor s\rfloor-3\ge s-4\) implies
\((1-\delta)^{h-3}\le(1-\delta)^{s-4}\).
If \(\delta=1\), the coupling already gives agreement with
probability one and \(s-4>0\), so the asserted upper bound is zero.
\end{proof}

In the following lemma, ``as above'' retains all the specification
hypotheses of Lemma~\ref{l:stopped}: consistency, full support on
nonempty finite site alphabets, range one in \(\ell^\infty\), a
uniform lower bound \(\alpha\in(0,1]\) on fully conditional
one-site probabilities, and (W) at every translate.

\begin{lemma}\label{l:hamming}
Assume (W) and \(\alpha\in(0,1]\) as above, and put
\(d_0=\min\{1/2,p_0\alpha^{16l_0}\}\).
Choose an integer \(n\) satisfying
\[
 n\ge44,\quad n^8\ge M_\star,\quad
 n^2\ge\max\{50000,6(l_0+1)\},\quad
 n^{12}d_0^{20}\ge50000\,2^{20}\,20!.
\]
Such an integer exists. Set \(M=n^8\).
Every translated \(C_M(a)\), with any two admissible boundaries
agreeing off one exterior site, admits a full-law coupling with
expected Hamming distance at most \(M/1000\).
One sufficient bound is
\[
 25n^6+25n^{16}(1-d_0)^{n-3}\le M/1000.
\]
\end{lemma}

\begin{proof}
The number \(d_0\) is strictly positive and at most \(1/2\).
All four requirements on \(n\) are lower bounds by finite constants;
for example, an integer at least
\[
 \max\left\{44,\ M_\star^{1/8},\
       \sqrt{\max\{50000,6(l_0+1)\}},\
       \left(50000\,2^{20}\,20!\,d_0^{-20}\right)^{1/12}\right\}
\]
satisfies them.  This proves existence before any torus periods or
boundary conditions are chosen.

We first prove a finite tail inequality.  For \(0<d\le1/2\) and
integer \(n\ge44\), put \(r=n-3\).  Then
\((1-d)^{-1}\ge1+d\), and the binomial theorem gives
\[
 \begin{split}
 (1-d)^{-r}
 &\ge(1+d)^r\ge\binom r{20}d^{20}\\
 &\ge\frac{(r-19)^{20}}{20!}d^{20}
   =\frac{(n-22)^{20}}{20!}d^{20}
   \ge\frac{n^{20}d^{20}}{2^{20}20!}.
 \end{split}
\]
Here \(r\ge41\), and each of the twenty factors in
\(r(r-1)\cdots(r-19)\) is at least \(r-19>0\).
Taking reciprocals proves
\[
 (1-d)^{n-3}\le\frac{2^{20}20!}{d^{20}n^{20}}.
\]

Fix a translated square with radius \(M=n^8\).  Its shell parameter
is \(s=M^{1/4}=n^2\), which is at least both \(16\) and
\(6(l_0+1)\).  Thus \(M\ge M_{\mathrm g}(l_0)\), and the other
size assumption gives \(M\ge M_\star\).
Apply Lemma~\ref{l:stopped} with \(h=n\); its index restriction
holds because \(4\le n\le n^2=\lfloor s\rfloor\).
If the changed exterior site is farther than distance one, use the
identical coupling and there is nothing to prove.  Otherwise the
lemma gives a full-law coupling agreeing outside the near region
\(\{x:|x-y|_\infty\le n^3\}\), except on an event of probability
at most \((1-\delta)^{n-3}\le(1-d_0)^{n-3}\).
The near region contains at most
\((2n^3+1)^2\le25n^6\) sites, and the square has
\((2n^8+1)^2\le25n^{16}\) sites.  Consequently the expected
Hamming cost of this same coupling is at most
\[
 25n^6+25n^{16}(1-d_0)^{n-3}.
\]
The condition \(n^2\ge50000\) gives \(25n^6\le n^8/2000\).
Using the tail inequality and the last requirement on \(n\) gives
\[
 25n^{16}(1-d_0)^{n-3}
   \le\frac{25\,2^{20}20!}{d_0^{20}n^4}
   \le\frac{n^8}{2000}.
\]
Adding proves the asserted bound \(M/1000\), uniformly over every
translate and boundary pair.  The truncation in the definition of
\(d_0\) also covers \(\delta=1\) without division by \(1-\delta\).
\end{proof}

In the following criterion the square radius and both torus periods
are integers.  The torus \(\ell^\infty\) distance is the minimum
over coordinate lifts.

\begin{theorem}\label{t:periodic}
Let \(M\ge1\), \(L,N\ge2M+3\), and let \(\mu>0\) be a
probability on a finite product over \(\Lambda_{L,N}\) whose
conditional law on any set depends only on letters in its exterior
\(\ell^\infty\) one-neighbourhood.
For every translated square \(Q\) of side \(2M+1\), suppose that:
the single-site heat bath of every conditional law on \(Q\) has Poincare gap at least
\(g>0\); and any two conditional block laws whose outside
configurations agree off one site admit a coupling with expected
Hamming distance in \(Q\) at most \(K\ge0\), whenever that site is
outside and at \(\ell^\infty\) distance one from \(Q\).
Put \(T_M=(2M+1)^2\), \(S_M=8M+8\), and suppose \(S_MK<T_M\).
Then
\[
 \sum_{Q}(I-E_Q)\text{ has Poincare gap at least }T_M-S_MK,\qquad
 \lambda(\mu)\ge g\left(1-\frac{S_MK}{T_M}\right).
\]
The sum is over all torus translates, each at rate one.
In particular, \(K\le M/1000\) implies
\(\lambda(\mu)\ge(499/500)g\).
\end{theorem}

\begin{proof}
For a centre label \(a\in\Lambda_{L,N}\), let
\(Q_a=a+\{-M,\ldots,M\}^2\), with coordinates reduced modulo
the periods.  There are \(A=LN\) labels.  We count these labels
without making any assumption on symmetries of \(\mu\).

Fix a site \(x\).  The centre labels for which \(x\in Q_a\) are
exactly \(a=x-z\), where \(z\in\{-M,\ldots,M\}^2\).
Reduction is injective on this offset square, so their number is
\(T_M=(2M+1)^2\).  Likewise, the labels with
\(\operatorname{dist}_\infty(x,Q_a)\le1\), using cyclic coordinate
distances on the torus, are exactly \(a=x-z\) with
\(z\in\{-M-1,\ldots,M+1\}^2\).  This follows by adding an
offset in \(\{-1,0,1\}^2\) to an offset in the smaller square;
their sum is precisely the larger square.  Its reduction is still
injective: a coordinate difference between two of its points has
absolute value at most \(2M+2\), strictly less than either relevant
period.  Hence exactly
\[
 (2M+3)^2-(2M+1)^2=8M+8=S_M
\]
labels exclude \(x\) but have it at distance one.  This count
includes the four corner offsets, which may have no interaction with
\(Q_a\) for a nearest-neighbour specification.  In particular, the
maximal site multiplicity of the labelled block family is \(T_M\).

Set \(P=A^{-1}\sum_aE_{Q_a}\).  Consider two input configurations
differing only at site \(x\), and choose the same uniformly random
block label in both components.  If \(x\in Q_a\), their exterior
configurations agree, so their conditional block laws are equal.
Use an identical sample on \(Q_a\); the output pair then agrees
everywhere.  If \(x\notin Q_a\) but its distance from \(Q_a\)
is one, use the assumed coupling on \(Q_a\).  The original
disagreement remains at \(x\), and its expected number of additional
disagreements is at most \(K\); the total is at most \(1+K\).
For every remaining label, the range-one assumption makes the two
block laws equal.  Identical resampling then leaves just the one
disagreement at \(x\).  The resulting mixture is a coupling of
the two \(P\)-output laws with
\[
 \mathbb E d_H
    \le\frac{S_M(1+K)+(A-T_M-S_M)}A
    =1-\frac{T_M-S_MK}A.
\]
By assumption \(c=T_M-S_MK>0\), and
\(c\le T_M\le A\) since \(K\ge0\).
Lemma~\ref{l:contraction} proves a block Poincare gap at least \(c\).
Applying Lemma~\ref{l:comparison} with the conditional block gap
\(g\) and multiplicity \(t_{\max}=T_M\) yields
\[
 \lambda(\mu)\ge\frac{cg}{T_M}
             =g\left(1-\frac{S_MK}{T_M}\right).
\]
These are the two claimed bounds with unnormalized block rates.
They also hold on one-point product spaces by the Poincare convention.
Finally, if \(K\le M/1000\), then
\[
 \frac{S_MK}{T_M}
 \le\frac{8M(M+1)}{1000(2M+1)^2}
 =\frac1{500}\left(1-\frac1{(2M+1)^2}\right)
 <\frac1{500}.
\]
In particular the strict block-contraction assumption is satisfied,
and the final stated lower bound follows.
\end{proof}



\begin{theorem}\label{t:mainW}
If the plane frame family satisfies (W), there exist \(c>0\) and
an integer \(L_0\ge3\) such that
\(\ker H_{L,N}=\mathbb C B_{L,N}\) and every nonzero eigenvalue is
at least \(c\), for all integers \(L,N\ge L_0\).
\end{theorem}

\begin{proof}
Fix constants \((l_0,p_0,M_\star)\) witnessing (W).
Lemma~\ref{l:planeenergy} supplies a consistent positive,
nearest-neighbour plane family with a common one-site bound
\(\alpha=12^{-1}r_b^4\in(0,1]\).  Choose \(n\) by
Lemma~\ref{l:hamming}, once and for all, and put
\[
 M=n^8,\qquad L_0=2M+3,\qquad
 g_M=r_b^{\,2(2M+1)(2M)}>0.
\]
These choices precede the torus periods, colouring, block placement
and exterior fibre.

Let \(L,N\ge L_0\).  In particular \(L,N\ge3\), so
Lemma~\ref{l:colour} provides the proper torus three-colouring
\[
 \kappa(i,j)=a_L(i)+a_N(j)\pmod3,\qquad
 a_m(i)=\begin{cases}
 2,&m\text{ odd and }i=m-1,\\
 i\bmod2,&\text{otherwise},
 \end{cases}
 \quad 0\le i<m.
\]
Consider
\(\mu=\mu^{\mathrm{af}}_{\kappa,L,N}\) from
Definition~\ref{d:framemodel}.  This measure has full support and
nearest-neighbour conditional dependence by
Lemma~\ref{l:planeenergy}.  The periodic pullback
\(\bar\kappa(z)=\kappa(p(z))\), where
\(p:\mathbb Z^2\to\Lambda_{L,N}\) is reduction, is a proper plane
colouring.  Indeed each plane nearest-neighbour edge reduces to a
torus edge between distinct vertices since \(L,N\ge3\).

For every labelled square \(Q\), choose its lift
\(Q'=a+\{-M,\ldots,M\}^2\) to the plane, and put
\(Q'^+=a+\{-M-1,\ldots,M+1\}^2\).
Here \(M=n^8\ge1\) and \(L,N\ge2M+3\), precisely the size
hypotheses of Lemma~\ref{l:boxlift}.  It gives injectivity of
\(p\) on \(Q'^+\) and the exact internal and crossing-bond
correspondences.  For an exterior torus configuration \(\sigma\),
the local sign change in that lemma is
\[
 \varepsilon_z=(-1)^{z_1+z_2},\qquad
 n_z=\varepsilon_zs_{p(z)}\quad(z\in Q'),\qquad
 \eta_z=\varepsilon_z\sigma_{p(z)}
                                  \quad(z\in Q'^+\setminus Q').
\]
Injectivity ensures \(p(z)\notin Q\) for
\(z\in Q'^+\setminus Q'\), so the displayed boundary values
are defined.  Only nearest-neighbour exterior letters enter the law.
Lemma~\ref{l:boxlift} identifies the conditional torus law under
this map with \(\mu^{+,\eta}_{\bar\kappa,Q'}\).
Antipodality of each alphabet
ensures that the sign change is a bijection using the same site
alphabets.  This bijection intertwines single-site updates and
preserves Hamming distances.  It also preserves variances and
Dirichlet forms, by the change of variables in the finite sums.
The number of internal edges of \(Q'\) is
\(2(2M+1)(2M)\): there are \(2M\) bonds in each of its
\(2M+1\) horizontal rows and the same number vertically.
Lemma~\ref{l:planeenergy}, equivalently the density comparison in
Lemma~\ref{l:density}, therefore gives conditional single-site
Poincare gap at least \(g_M\), uniformly over all such fibres.

Now take two torus exterior configurations for \(Q\) differing at
just one site \(y\) at \(\ell^\infty\) distance one.
The image \(p(Q'^+)\) is the full torus one-neighbourhood of \(Q\):
every one-step coordinate displacement lifts to
\(\{-1,0,1\}^2\), and
\(Q'^+=Q'+\{-1,0,1\}^2\).
Thus the one-neighbourhood injection in Lemma~\ref{l:boxlift}
gives an existing, unique exterior lift of \(y\).
Under the same local
sign change the two plane boundary assignments differ at most at
that lift.  Only nearest-neighbour exterior letters are needed; if
necessary, extend the other exterior letters identically and
arbitrarily in their site alphabets.  For a corner site that has no
crossing nearest-neighbour bond the two conditional laws are already
equal.  In all other cases Lemma~\ref{l:hamming}, applied to this
plane colouring and placement, provides a coupling of expected
Hamming cost at most \(K=M/1000\).  The inverse coordinatewise
bijection carries it to a coupling of the two conditional torus laws
with exactly the same Hamming cost.

All hypotheses of Theorem~\ref{t:periodic} now hold, giving
\[
 \lambda(\mu^{\mathrm{af}}_{\kappa,L,N})
                        \ge\frac{499}{500}g_M.
\]
Use the constant defined and proved positive in Lemma~\ref{l:star}:
\(g_\star\) is the smallest positive eigenvalue of
\(h_\star=\sum_{j=1}^4P_{\{0,j\}}\) on \(V_4^{\otimes5}\).
It is independent of the torus periods.  Theorem~\ref{t:QD} applies
to every \(L,N\ge3\) and every proper torus three-colouring using
the frames of Definition~\ref{d:framemodel}; these hypotheses hold
for the periods and colouring just chosen.  Its precise operator
inequality is
\[
 H_{L,N}\big|_{B_{L,N}^{\perp}}
       \ge\frac{g_\star}{2}
               \lambda(\mu^{\mathrm{af}}_{\kappa,L,N})I.
\]
The factor \(1/2\) uses Lemma~\ref{l:star}:
\(\sum_xH_x=2H_{L,N}\), where \(H_x=\sum_{e\ni x}P_e\),
with one heat-bath update per site.  Consequently every eigenvalue on
\(B_{L,N}^{\perp}\) is at least
\[
 c=\frac{g_\star}{2}\,\frac{499}{500}g_M
                         =\frac{499}{1000}g_\star g_M>0.
\]
Proposition~\ref{p:toruskernel}, whose hypothesis is only
\(L,N\ge3\), supplies both \(B_{L,N}\ne0\) and
\(\ker H_{L,N}=\mathbb C B_{L,N}\).  Thus \(B_{L,N}^{\perp}\)
is the direct sum of the eigenspaces for nonzero eigenvalues.
The constants \(c,L_0\) are independent of \(L,N\).
The sign change was defined separately on each lifted square, so this
argument applies to both even and odd periods.
\end{proof}


We now prove the two conclusions of Theorem~\ref{t:main}, stated in
the result section: nonvanishing of \(B_{L,N}\) and the kernel
identity for every \(L,N\ge3\), and the existence of one
\(c>0\) and integer \(L_0\ge3\) bounding every nonzero eigenvalue
below by \(c\) for all integers \(L,N\ge L_0\).

\begin{proof}[Proof of Theorem~\ref{t:main}]
Proposition~\ref{p:toruskernel} proves \(B_{L,N}\ne0\) and
\(\ker H_{L,N}=\mathbb C B_{L,N}\) for every \(L,N\ge3\).
For the uniform positive gap, Theorem~\ref{t:inputs} gives (H1),
(H2) with \(C_V=10\) and
\(0<q=1119948287353/1120437614225<1\), and (H3\(_h\)), all at
the same pair \((c_{\rm s},c_{\rm t})=(961/1000,187/200)\).
Theorem~\ref{t:rec} thus gives \(A\ge2\), \(0<\theta<1\) and
\(\mathsf P(A,\theta)\).
Theorem~\ref{t:cutMOS} supplies
\(\mathsf W(64,p_0,M_\star)\) with common
\(p_0\in(0,1]\) and positive integer \(M_\star\), valid for
every proper plane colouring and every translate.
Theorem~\ref{t:mainW} then gives \(c>0\) and an integer
\(L_0\ge3\) such that all nonzero eigenvalues are at least \(c\)
whenever \(L,N\ge L_0\), as required.
\end{proof}
